\begingroup
\renewcommand{\arraystretch}{1.45}
\InterpolationArticle{20}
Betrachten wir zunächst den allgemeinen Fall von Artikel 10, in dem \(X\) die Form
\[
\begin{aligned}
&\alpha+\alpha'\cos x+\alpha''\cos2x+\alpha'''\cos3x+\mathinner{.\,.\,.}+\alpha^{(m)}\cos mx\\[5pt]
&\quad+\beta'\sin x+\beta''\sin2x+\beta'''\sin3x+\mathinner{.\,.\,.}+\beta^{(m)}\sin mx
\end{aligned}
\]
\noindent hat und \(\mu=2m+1\) ist. Hier gilt daher
\[
\begin{aligned}
&\sin\frac12(t-b)\sin\frac12(t-c)\sin\frac12(t-d)\mathinner{.\,.\,.}=\frac{\sin\frac12\mu(t-a)}{2^{\mu-1}\sin\frac12(t-a)}\\[5pt]
&=\frac{1}{2^{\mu-1}}\bigl(1+2\cos(t-a)+2\cos2(t-a)\\[5pt]&\qquad{}+2\cos3(t-a)+\mathinner{.\,.\,.}+2\cos m(t-a)\bigr)
\end{aligned}
\]
\noindent und somit, wenn \(a\) für \(t\) eingesetzt wird,
\[
\begin{aligned}
&\sin\frac12(a-b)\sin\frac12(a-c)\sin\frac12(a-d)\mathinner{.\,.\,.}=\frac{\mu}{2^{\mu-1}}
\end{aligned}
\]
\SourceLeaf{305}
In der Formel von Artikel 10 für \(T\) wird also der Koeffizient von \(A\)
\[
\begin{aligned}
&=\frac{1}{\mu}(1+2\cos(t-a)+2\cos2(t-a)\\[5pt]&\qquad{}+2\cos3(t-a)+\mathinner{.\,.\,.}+2\cos m(t-a))
\end{aligned}
\]
Ebenso gilt
\[
\begin{aligned}
&\sin\frac12(t-a)\sin\frac12(t-c)\sin\frac12(t-d)\mathinner{.\,.\,.}
=-\frac{\sin\frac12\mu(t-b)}{2^{\mu-1}\sin\frac12(t-b)}\\[5pt]
&=-\frac{1}{2^{\mu-1}}(1+2\cos(t-b)+2\cos2(t-b)\\[5pt]&\qquad{}+2\cos3(t-b)+\mathinner{.\,.\,.}+2\cos m(t-b))
\end{aligned}
\]
und
\[
\begin{aligned}
&\sin\frac12(b-a)\sin\frac12(b-c)\sin\frac12(b-d)\mathinner{.\,.\,.}=-\frac{\mu}{2^{\mu-1}}
\end{aligned}
\]
Somit wird der Koeffizient von \(B\) in der Formel für \(T\)
\[
\begin{aligned}
&=\frac{1}{\mu}(1+2\cos(t-b)+2\cos2(t-b)\\[5pt]&\qquad{}+2\cos3(t-b)+\mathinner{.\,.\,.}+2\cos m(t-b))
\end{aligned}
\]
Die Koeffizienten von \(C,D\) usw. werden ganz ähnlich sein; daraus folgt schließlich,
dass
\[
\begin{aligned}
&T=\frac{1}{\mu}(A+B+C+D+\mathinner{.\,.\,.})\\[5pt]
&+\frac{2}{\mu}(A\cos a+B\cos b+C\cos c+D\cos d+\mathinner{.\,.\,.})\cos t\\[5pt]
&+\frac{2}{\mu}(A\sin a+B\sin b+C\sin c+D\sin d+\mathinner{.\,.\,.})\sin t\\[5pt]
&+\frac{2}{\mu}(A\cos2a+B\cos2b+C\cos2c+D\cos2d+\mathinner{.\,.\,.})\cos2t\\[5pt]
&+\frac{2}{\mu}(A\sin2a+B\sin2b+C\sin2c+D\sin2d+\mathinner{.\,.\,.})\sin2t\\[5pt]
&+\mathrm{usw.}\\[5pt]
&+\frac{2}{\mu}(A\cos ma+B\cos mb+C\cos mc+D\cos md+\mathinner{.\,.\,.})\cos mt\\[5pt]
&+\frac{2}{\mu}(A\sin ma+B\sin mb+C\sin mc+D\sin md+\mathinner{.\,.\,.})\sin mt
\end{aligned}
\]
Da diese Formel mit
\[
\begin{aligned}
&\alpha+\alpha'\cos t+\beta'\sin t+\alpha''\cos2t+\beta''\sin2t+\mathrm{usw.}\\[5pt]&\qquad{}+\alpha^{(m)}\cos mt+\beta^{(m)}\sin mt
\end{aligned}
\]
identisch sein muss, ergeben sich daraus unmittelbar die Werte der Koeffizienten \(\alpha,\alpha',\beta',\alpha'',\beta''\)
usw.
\InterpolationArticle{21}
Bricht die Reihe für \(X\) nicht mit den Gliedern \(\cos mx\) und \(\sin mx\) ab, ob
sie nun ins Unendliche fortläuft oder endlich ist, so wird der im vorigen Artikel für \(T\) gefundene Wert unvoll\SourceLeaf{306}%
ständig sein; ebenso werden die dort angegebenen Werte der einzelnen Koeffizienten \(\alpha,\alpha',\beta',\alpha'',\beta''\) usw.
einer Korrektur bedürfen. Diese lässt sich durch die folgende Methode bequemer bestimmen als durch die Summation von Artikel 2. Zunächst bemerken wir, dass
\[
\begin{aligned}
&\cos na+\cos nb+\cos nc+\cos nd+\mathinner{.\,.\,.}=\mu\cos na\\[5pt]
&\sin na+\sin nb+\sin nc+\sin nd+\mathinner{.\,.\,.}=\mu\sin na
\end{aligned}
\]
gilt, wenn \(n\) eine durch \(\mu\) teilbare ganze Zahl ist; dagegen gilt
\[
\begin{aligned}
&\cos na+\cos nb+\cos nc+\cos nd+\mathinner{.\,.\,.}=0\\[5pt]
&\sin na+\sin nb+\sin nc+\sin nd+\mathinner{.\,.\,.}=0
\end{aligned}
\]
wenn \(n\) eine durch \(\mu\) nicht teilbare ganze Zahl ist. Im ersten Fall ist die Sache selbstverständlich;
im zweiten sei die erste Summe \(=P\), die zweite \(=Q\), woraus leicht folgt, dass
\[
\begin{aligned}
&P\cos n(b-a)-Q\sin n(b-a)=P\\[5pt]
&P\sin n(b-a)+Q\cos n(b-a)=Q
\end{aligned}
\]
Multipliziert man die erste Gleichung mit \(\cos n(b-a)-1\), die zweite mit \(\sin n(b-a)\),
so ergibt sich durch Addition und Streichung der sich gegenseitig aufhebenden Glieder
\[
\begin{aligned}
&2P(1-\cos n(b-a))=0
\end{aligned}
\]
Ebenso erhält man, wenn die erste Gleichung mit \(\sin n(b-a)\), die zweite mit
\(1-\cos n(b-a)\) multipliziert wird, durch Addition
\[
\begin{aligned}
&2Q(1-\cos n(b-a))=0
\end{aligned}
\]
Im ersten Fall nun (in dem \(n\) durch \(\mu\) teilbar ist und daher \(\cos n(b-a)=1\) gilt)
sind diese Gleichungen Identitäten; im zweiten dagegen (in dem \(n\) nicht durch \(\mu\) teilbar ist und somit \(\cos n(b-a)=\cos\frac{n}{\mu}\times360^{0}\) nicht \(=1\) sein kann) können sie nur gelten, wenn \(P=0\) und \(Q=0\) gilt.

\Needspace{30\baselineskip}
Setzen wir nun, dass auf die Glieder \(\alpha^m\cos mx+\beta^m\sin mx\) im Ausdruck der Funktion \(X\)
\[
\begin{aligned}
&\phantom{{}+{}}\alpha^{m+1}\cos(m+1)x+\alpha^{m+2}\cos(m+2)x+\mathrm{usw.}\\[5pt]
&+\beta^{m+1}\sin(m+1)x+\beta^{m+2}\sin(m+2)x+\mathrm{usw.}
\end{aligned}
\]
folgen und dass der im vorigen Artikel unter der Voraussetzung, diese Teile fehlten, für \(T\) gefundene (unvollständige)
Wert
\SourceLeaf{307}
\[
\begin{aligned}
&\gamma+\gamma'\cos t+\gamma''\cos2t+\gamma'''\cos3t+\mathinner{.\,.\,.}+\gamma^m\cos mt\\[5pt]
&\phantom{\gamma}\;+\delta'\sin t+\delta''\sin2t+\delta'''\sin3t+\mathinner{.\,.\,.}+\delta^m\sin mt
\end{aligned}
\]
ist, sodass \(\gamma=\frac{1}{\mu}(A+B+C+D+\mathrm{usw.})\) usw. gilt. Für \(A,B,C,D\)
usw. setze man ihre Werte ein, nämlich
\[
\begin{aligned}
&A=\alpha+\alpha'\cos a+\alpha''\cos2a+\mathinner{.\,.\,.}+\alpha^m\cos ma\\[5pt]&\qquad{}+\alpha^{m+1}\cos(m+1)a\mathinner{.\,.\,.}\\[5pt]
&+\beta'\sin a+\beta''\sin2a+\mathinner{.\,.\,.}+\beta^m\sin ma\\[5pt]&\qquad{}+\beta^{m+1}\sin(m+1)a\mathinner{.\,.\,.}\\[5pt]
&\mathrm{usw.}
\end{aligned}
\]
und es ergibt sich durch die eben angegebene Summation
\[
\begin{aligned}
&\gamma=\alpha+\alpha^{\mu}\cos\mu a+\alpha^{2\mu}\cos2\mu a+\alpha^{3\mu}\cos3\mu a+\mathrm{usw.}\\[5pt]
&+\beta^{\mu}\sin\mu a+\beta^{2\mu}\sin2\mu a+\beta^{3\mu}\sin3\mu a+\mathrm{usw.}\\[5pt]
&\gamma'=\alpha'+(\alpha^{\mu-1}+\alpha^{\mu+1})\cos\mu a\\[5pt]&\qquad{}+(\alpha^{2\mu-1}+\alpha^{2\mu+1})\cos2\mu a+\mathrm{usw.}\\[5pt]
&+(\beta^{\mu-1}+\beta^{\mu+1})\sin\mu a\\[5pt]&\qquad{}+(\beta^{2\mu-1}+\beta^{2\mu+1})\sin2\mu a+\mathrm{usw.}\\[5pt]
&\delta'=\beta'-(\beta^{\mu-1}-\beta^{\mu+1})\cos\mu a\\[5pt]&\qquad{}-(\beta^{2\mu-1}-\beta^{2\mu+1})\cos2\mu a-\mathrm{usw.}\\[5pt]
&-(\alpha^{\mu-1}-\alpha^{\mu+1})\sin\mu a\\[5pt]&\qquad{}-(\alpha^{2\mu-1}-\alpha^{2\mu+1})\sin2\mu a-\mathrm{usw.}\\[5pt]
&\gamma''=\alpha''+(\alpha^{\mu-2}+\alpha^{\mu+2})\cos\mu a\\[5pt]&\qquad{}+(\alpha^{2\mu-2}+\alpha^{2\mu+2})\cos2\mu a+\mathrm{usw.}\\[5pt]
&+(\beta^{\mu-2}+\beta^{\mu+2})\sin\mu a\\[5pt]&\qquad{}+(\beta^{2\mu-2}+\beta^{2\mu+2})\sin2\mu a+\mathrm{usw.}\\[5pt]
&\delta''=\beta''-(\beta^{\mu-2}-\beta^{\mu+2})\cos\mu a\\[5pt]&\qquad{}-(\beta^{2\mu-2}-\beta^{2\mu+2})\cos2\mu a-\mathrm{usw.}\\[5pt]
&-(\alpha^{\mu-2}-\alpha^{\mu+2})\sin\mu a\\[5pt]&\qquad{}-(\alpha^{2\mu-2}-\alpha^{2\mu+2})\sin2\mu a-\mathrm{usw.}
\end{aligned}
\]
usw., bis zu
\[
\begin{aligned}
&\gamma^m=\alpha^m+(\alpha^{\mu-m}+\alpha^{\mu+m})\cos\mu a\\[5pt]&\qquad{}+(\alpha^{2\mu-m}+\alpha^{2\mu+m})\cos2\mu a+\mathrm{usw.}\\[5pt]
&+(\beta^{\mu-m}+\beta^{\mu+m})\sin\mu a\\[5pt]&\qquad{}+(\beta^{2\mu-m}+\beta^{2\mu+m})\sin2\mu a+\mathrm{usw.}\\[5pt]
&\delta^m=\beta^m-(\beta^{\mu-m}-\beta^{\mu+m})\cos\mu a\\[5pt]&\qquad{}-(\beta^{2\mu-m}-\beta^{2\mu+m})\cos2\mu a-\mathrm{usw.}\\[5pt]
&-(\alpha^{\mu-m}-\alpha^{\mu+m})\sin\mu a\\[5pt]&\qquad{}-(\alpha^{2\mu-m}-\alpha^{2\mu+m})\sin2\mu a-\mathrm{usw.}
\end{aligned}
\]
Konvergiert daher die Reihe so stark, dass \(\alpha^{m+1},\beta^{m+1},\alpha^{m+2},\beta^{m+2}\) usw. vernachlässigt werden
können, so darf man die Werte \(\gamma,\gamma',\delta'\) usw. für die wahren Werte \(\alpha,\alpha',\beta'\) usw. nehmen.
\endgroup
